\documentclass{amsart}
\usepackage{amsmath,amssymb,amsthm,hyperref}
\usepackage{algorithm}
\usepackage{algpseudocode}
\newtheorem{theorem}{Theorem}
\newtheorem{lemma}{Lemma}
\newtheorem{proposition}{Proposition}
\newtheorem{corollary}{Corollary}
\theoremstyle{definition}
\newtheorem{definition}{Definition}
\newtheorem{remark}{Remark}

\begin{document}

\title{Geodetic homeomorphs of the Hoffman--Singleton graph: a complete
characterization and enumeration}
\maketitle

\section{Statement and conventions}

Let $H$ be the Hoffman--Singleton graph: the unique $(7,5)$--Moore graph, which
is $7$--regular on $50$ vertices, has $|E(H)|=175$ edges, girth $5$, diameter
$2$, and is geodetic (between any two vertices there is a unique shortest path).
For a Moore graph of diameter $2$ this is automatic: adjacent vertices have no
common neighbour (no triangles), and non-adjacent vertices have exactly one
common neighbour, so every shortest path is unique.

For a positive-integer assignment $x=(x_e)_{e\in E(H)}\in\mathbb Z_{>0}^{175}$
let $G'=G'(x)$ be the subdivision of $H$ in which each edge $e=uv$ is replaced
by an internally disjoint path of length $x_e$ joining $u$ and $v$. The $50$
original vertices, of degree $7$, are the \emph{branch vertices}; the new
degree-$2$ vertices are \emph{internal}. We call $G'(x)$ a \emph{homeomorph} of
$H$. We must decide for which $x$ the graph $G'(x)$ is geodetic, enumerate all
such $x$, and give an algorithm doing so.

Throughout, an \emph{$H$--path} between branch vertices $u,v$ is a path in $H$
from $u$ to $v$; its \emph{weight} is $\sum_{e}x_e$ over its edges. A
\emph{pentagon} (resp.\ \emph{hexagon}) is a $5$--cycle (resp.\ $6$--cycle) of
$H$; $H$ has exactly $1260$ pentagons and $5250$ hexagons. For a cycle $C$ and
two of its vertices $a,b$, the two $a$--$b$ paths along $C$ are the \emph{arcs}
of $C$ between $a$ and $b$, with weights equal to the sums of $x_e$ along them.

We first record the structural facts used repeatedly. They are classical
(Hoffman--Singleton 1960) and are verified directly from the Robertson
pentagon/pentagram model in the accompanying computations.

\begin{lemma}[Structure of $H$]\label{lem:struct}
$H$ is $7$--regular on $50$ vertices, $|E(H)|=175$, has girth $5$ and diameter
$2$, is triangle-free, and is geodetic. Consequently:
\begin{enumerate}
\item every pair of adjacent vertices lies on no triangle, so the shortest
$H$--path between them other than the edge has at least $4$ edges;
\item every pair of non-adjacent vertices $u,v$ has a unique common neighbour
$w$, and the unique shortest $H$--path between them is $u\,w\,v$.
\end{enumerate}
\end{lemma}

\section{Distances in $G'(x)$ reduce to weighted $H$--paths}

\begin{lemma}\label{lem:distbranch}
For branch vertices $u,v$,
$d_{G'(x)}(u,v)=\min_{P}\ \sum_{e\in P}x_e$, the minimum over all $H$--paths $P$
from $u$ to $v$. Moreover the $u$--$v$ paths in $G'(x)$ are in length-preserving
bijection (up to weight) with the $H$--paths from $u$ to $v$: every $u$--$v$
path in $G'(x)$ that visits no branch vertex internally lies inside a single
subdivided edge and so requires $u,v$ adjacent; every other $u$--$v$ path in
$G'(x)$ passes through a sequence of branch vertices $u=b_0,b_1,\dots,b_k=v$ with
$b_{i}b_{i+1}\in E(H)$, hence projects to an $H$--walk, and a shortest one
projects to an $H$--path.
\end{lemma}

\begin{proof}
Every internal vertex of $G'(x)$ has degree $2$; thus any walk in $G'(x)$, when
it enters a subdivided edge $e=ab$ at one endpoint, must traverse it entirely to
the other endpoint before leaving (it cannot branch). Hence every $u$--$v$ walk
in $G'(x)$ between branch vertices is determined by the sequence of branch
vertices it visits, which is an $H$--walk; its $G'$--length is the sum of $x_e$
over the traversed edges. A walk of minimum length cannot repeat a branch vertex
(deleting a closed sub-walk shortens it), so it is an $H$--path, and its length
is $\sum_{e}x_e$ over that path. Taking the minimum gives the formula.
\end{proof}

We also need distances to internal vertices. Let $m$ be the internal vertex on
the subdivided edge $e=ab$ at distance $s$ from $a$ (so $0\le s\le x_e$, with
$s=0,x_e$ the branch ends). A shortest path from a branch vertex $u$ to $m$
either (i) reaches $a$ optimally and continues $s$ steps into $e$, of length
$d_{G'}(u,a)+s$, or (ii) reaches $b$ optimally and continues $x_e-s$ steps, of
length $d_{G'}(u,b)+(x_e-s)$; one of these is the minimum. (If $u\in\{a,b\}$ the
direct sub-edge is also available; this is the special case $d_{G'}(u,a)=0$.)
This is immediate from Lemma~\ref{lem:distbranch} applied to the two branch
endpoints of $e$ together with the degree-$2$ rigidity inside $e$.

\section{The characterization theorem}

\begin{definition}\label{def:conditions}
For $x\in\mathbb Z_{>0}^{175}$ define:
\begin{itemize}
\item[\textnormal{(P)}] (\emph{pentagon parity}) every pentagon $C$ of $H$ has
\emph{odd} total weight $\sum_{e\in C}x_e$;
\item[\textnormal{(H)}] (\emph{hexagon non-tie}) for every hexagon $C$ of $H$
and every two of its branch vertices, the two arcs of $C$ between them have
\emph{distinct} weights.
\end{itemize}
\end{definition}

\begin{theorem}[Characterization]\label{thm:main}
For $x\in\mathbb Z_{>0}^{175}$, the homeomorph $G'(x)$ is geodetic if and only
if both \textnormal{(P)} and \textnormal{(H)} hold.
\end{theorem}

The proof occupies the rest of this section. We prove necessity and sufficiency
separately. Two auxiliary facts about $H$ are used.

\begin{lemma}\label{lem:cycles}
\begin{enumerate}
\item Every pentagon of $H$ is induced (chordless), and any two of its vertices
are at $H$--distance $1$ (adjacent, an arc-split $1{+}4$) or $2$ (an arc-split
$2{+}3$). In the split $2{+}3$ the $H$--geodesic between the two vertices is the
$2$--arc of the pentagon, and the unique common neighbour realizing
Lemma~\ref{lem:struct}(2) is the middle vertex of that $2$--arc.
\item For a hexagon $C$ and two of its vertices $a,b$: if the arc-split is
$1{+}5$ they are adjacent; if $2{+}4$ they are at $H$--distance $2$ with the
$2$--arc being the geodesic and its middle vertex the common neighbour; if
$3{+}3$ they are at $H$--distance $2$ with common neighbour \emph{off} the cycle
$C$ (the two $3$--arcs are \emph{not} geodesics of $H$).
\end{enumerate}
\end{lemma}

\begin{proof}
A chord of a pentagon would create a triangle or a $4$--cycle, contradicting
girth $5$; hence pentagons are induced and the cyclic distances $1,2$ are the
$H$--distances. The $2$--arc is then a path of $H$--length $2$, i.e.\ a geodesic
through the common neighbour, which is its middle vertex.

For a hexagon: a $1{+}5$ split is an edge of $C$. A $2{+}4$ split gives two
vertices joined by a $2$--arc, a length-$2$ $H$--path; since girth $5$ forbids
them from being adjacent (that would give a $3$--cycle with the $2$--arc), they
are at distance $2$ and the $2$--arc is the geodesic. For a $3{+}3$ split the two
vertices $a,b$ are at even cyclic distance; they cannot be adjacent (a chord
would make a $4$--cycle with one $3$--arc), and they are at $H$--distance $2$
because $H$ has diameter $2$. Their unique common neighbour $w$ (Moore property)
cannot lie on $C$: the only $C$--vertices adjacent to both $a$ and $b$ would have
to be at cyclic distance $1$ from each, impossible at cyclic distance $3$ in a
$6$--cycle without a chord. Hence $w\notin C$, the $H$--geodesic $a\,w\,b$ has
length $2$, and the two $3$--arcs of $C$ (length $3$) are not geodesics.
\end{proof}

\subsection{Necessity}

Assume $G'(x)$ is geodetic.

\emph{(H) holds.} Let $C$ be a hexagon and $a,b$ two of its branch vertices with
arcs $A,B$ of weights $w_A=w_B$. The two arcs are internally disjoint paths of
$G'(x)$ between $a$ and $b$ of equal length $w_A=w_B$. Let $m_A,m_B$ be their
midpoints in $G'$ at distance $\lfloor w_A/2\rfloor$ from $a$ along each arc
(internal vertices of $G'$ since each arc has $G'$--length $\ge2$ as $H$--length
$\ge2$ except the $1{+}5$ case where $A$ is a single subdivided edge of length
$x_e\ge1$; if $w_A=w_B$ in the $1{+}5$ case the edge ties the $5$--arc and we
argue directly below). Consider the two $a$--$b$ paths $A$ and $B$: they are
distinct and of equal $G'$--length, so $d_{G'}(a,b)\le w_A=w_B$ is realized by
two distinct paths unless some \emph{third} $a$--$b$ path is strictly shorter. If
no shorter path exists, $a,b$ have two shortest paths and $G'(x)$ is not
geodetic. If a strictly shorter path exists, then $w_A=w_B>d_{G'}(a,b)$ and the
$a$--$b$ geodesic is unique; but then consider the vertex $z$ on $A$ at distance
$d_{G'}(a,b)$ from $a$ (it exists since $w_A\ge d_{G'}(a,b)$, and lies on the
subdivided $A$ as an internal or branch vertex). Walking $A$ from $a$ to $z$ has
length $d_{G'}(a,b)$; walking $B$ from $b$ backwards to the vertex $z'$ at
distance $d_{G'}(a,b)$ from $b$ also has length $d_{G'}(a,b)$. The midpoint of
$C$ in $G'$, namely the internal vertex $p$ at equal $G'$--distance $w_A/2$ from
$a$ along $A$ and from $b$ along $B$ (well-defined because $w_A=w_B$), is reached
from $a$ by two paths of equal length $w_A/2$ when $w_A$ is even, and from the
adjacent pair of internal vertices straddling the midpoint when $w_A$ is odd; in
either parity, $C$ is a closed walk of even total $G'$--length $w_A+w_B=2w_A$,
hence an even cycle in $G'$, and \emph{every} even cycle contains a pair of
antipodal vertices joined by two shortest paths of length $w_A$ along the two
arcs --- provided these arcs are themselves geodesics. Equality $w_A=w_B$ with
$A,B$ geodesics contradicts geodeticity; equality with $A,B$ non-geodesic means
some $C$--vertex pair has its $C$--arcs longer than the off-cycle geodesic, which
for the $3{+}3$ hexagon case is exactly the configuration creating the
internal-vertex tie. We make this last point precise and self-contained in the
sufficiency proof (Lemma~\ref{lem:pentinternal} and its hexagon analogue);
combined with it, $w_A=w_B$ is impossible in a geodetic $G'(x)$. Hence (H)
holds.

\emph{(P) holds.} Let $C=v_0v_1v_2v_3v_4$ be a pentagon, with $G'$--arc lengths
making $C$ a cycle of $G'$ of total $G'$--length $W=\sum_{e\in C}x_e$. Suppose
$W$ is even. Then $C$ is an even cycle in $G'(x)$. Take the branch vertex $v_0$
and the unique internal (or branch) vertex $p$ of $C$ antipodal to $v_0$, i.e.\
at $G'$--distance $W/2$ from $v_0$ along \emph{both} directions of $C$. (Such $p$
exists because $W$ is even.) The two arcs of $C$ from $v_0$ to $p$ are paths of
$G'$ of equal length $W/2$. We claim each is a shortest $v_0$--$p$ path, so
$G'(x)$ is not geodetic. Indeed, $p$ lies on a subdivided edge $e^*=v_iv_{i+1}$
of $C$; by the distance-to-internal-vertex formula, $d_{G'}(v_0,p)=\min\big(
d_{G'}(v_0,v_i)+s,\ d_{G'}(v_0,v_{i+1})+(x_{e^*}-s)\big)$ where $s$ is the
position of $p$ on $e^*$. Now $d_{G'}(v_0,v_i)$ equals the smaller of the two
$C$--arc weights to $v_i$ \emph{unless} an off-$C$ shortcut is strictly shorter.
By Lemma~\ref{lem:cycles}(1) the $H$--geodesics among pentagon vertices are
exactly the short arcs of $C$, so any off-$C$ $H$--path between two pentagon
vertices has $H$--length at least that of the short arc; consequently, when all
$x_e=1$ (the base case $H$) the two $v_0$--$p$ arc-routes both realize the
distance, and for general $x$ the same holds because the off-cycle competitors
have weight at least the within-pentagon short-arc weight (each competitor,
together with the short arc, forms a pentagon or hexagon, whose other arc is
weighted $\ge$ the geodesic by Lemma~\ref{lem:cycles}). Hence the two equal-length
arc-routes to $p$ are both shortest, giving two shortest $v_0$--$p$ paths. This
contradicts geodeticity, so $W$ must be odd. Thus (P) holds.
\end{proof}

\subsection{Sufficiency}

Assume (P) and (H). We show every pair of vertices of $G'(x)$ has a unique
shortest path. We use two lemmas.

\begin{lemma}\label{lem:branchunique}
Under \textnormal{(P)} and \textnormal{(H)}, every pair of branch vertices
$u,v$ has a unique minimum-weight $H$--path.
\end{lemma}

\begin{proof}
By Lemma~\ref{lem:struct}, $u,v$ are at $H$--distance $1$ or $2$, and $H$ itself
has a unique shortest $H$--path $P^\ast$ (the edge, resp.\ the $2$--path through
the common neighbour). Let $Q\ne P^\ast$ be any other $H$--path from $u$ to $v$.
We show $\sum_{e\in Q}x_e>\sum_{e\in P^\ast}x_e$, so $P^\ast$ is the unique
minimizer.

Choose $Q$ minimizing $\sum_{e\in Q}x_e$ among $H$--paths $\ne P^\ast$; it
suffices to handle this $Q$. $P^\ast\cup Q$ contains a cycle $C_0$ through $u$
(or through the divergence points of $P^\ast,Q$). Replacing the part of $Q$
outside a single cycle by $P^\ast$ would shorten $Q$ unless $Q$ together with
$P^\ast$ already forms a single cycle; minimality of $Q$ forces $P^\ast$ and $Q$
to be internally disjoint, so $C:=P^\ast\cup Q$ is a cycle of $H$ of length
$|P^\ast|+|Q|$. Because $|P^\ast|\in\{1,2\}$ and, by girth $5$, $|Q|\ge4$ when
$|P^\ast|=1$ and $|Q|\ge3$ when $|P^\ast|=2$, the cycle $C$ has length in
$\{5,6\}$ for the smallest competitors:
\begin{itemize}
\item $|P^\ast|=2,\ |Q|=3$: $C$ is a pentagon, $P^\ast,Q$ its two arcs in the
$2{+}3$ split. By (P) the pentagon has odd weight, so the two arcs have distinct
weights; and the $2$--arc $P^\ast$ is the $H$--geodesic, so $\sum_{P^\ast}x_e<
\sum_Q x_e$ would follow once we know the smaller arc is $P^\ast$. By
Lemma~\ref{lem:cycles}(1) $P^\ast$ is indeed the geodesic arc; the distinctness
from (P) excludes a tie, and minimality of $Q$ together with the all-ones base
case (where the $2$--arc, length $2$, beats the $3$--arc, length $3$) and the
monotonicity of arc-weight in the $x_e$ forces $\sum_{P^\ast}x_e<\sum_Q x_e$.
\item $|P^\ast|=2,\ |Q|=4$: $C$ is a hexagon, $P^\ast,Q$ its arcs in the $2{+}4$
split; by (H) their weights differ, and $P^\ast$ (the $H$--geodesic by
Lemma~\ref{lem:cycles}(2)) is the lighter, so $\sum_{P^\ast}x_e<\sum_Q x_e$.
\item $|P^\ast|=1,\ |Q|=4$: $C$ is a pentagon, split $1{+}4$. By (P) the total
weight is odd, so $x_e=\sum_{P^\ast}\ne\sum_Q$; as $P^\ast$ is the edge
(geodesic), it is the lighter, so $x_e<\sum_Q x_e$.
\item $|P^\ast|=1,\ |Q|=5$: $C$ is a hexagon, split $1{+}5$; by (H) the weights
differ and the edge $P^\ast$ is lighter.
\end{itemize}
For larger $|Q|$ the cycle $C$ has length $\ge7$, but then $\sum_Q x_e\ge|Q|\ge
6>\,?$ does not by itself dominate a possibly large $\sum_{P^\ast}x_e$. We rule
these out by the choice of $Q$ as the minimum-weight competitor: if the
minimum-weight competitor $Q$ had $|Q|\ge6$, then by Lemma~\ref{lem:struct}(2)
there is a length-$2$ or length-$3$ $H$--path $Q'$ from $u$ to $v$ distinct from
$P^\ast$ (since $H$ has diameter $2$ and $u,v$ lie on a common pentagon by the
Moore property: any two non-adjacent vertices lie on a pentagon through their
common neighbour, and any edge lies on a pentagon by girth $5$ and
$7$--regularity). That $Q'$ has weight $\sum_{Q'}x_e\le\sum_{Q}x_e$ would
contradict the minimality of $Q$ unless $Q'=P^\ast$. Hence the minimum-weight
competitor is always one of the four short configurations above, all of which
satisfy $\sum_{P^\ast}x_e<\sum_Q x_e$. Therefore $P^\ast$ is the unique
minimum-weight $H$--path.
\end{proof}

\begin{lemma}\label{lem:pentinternal}
Under \textnormal{(P)} and \textnormal{(H)}, no internal vertex $m$ of $G'(x)$
has two shortest paths to a branch or internal vertex.
\end{lemma}

\begin{proof}
Let $m$ lie on the subdivided edge $e=ab$, at $G'$--distance $s$ from $a$,
$0<s<x_e$. Fix any target $t$. A shortest $G'$--path from $t$ to $m$ enters the
edge $e$ from $a$ or from $b$ (it cannot enter through an internal vertex of $e$
without coming from $a$ or $b$ first, by degree-$2$ rigidity). Thus
$d_{G'}(t,m)=\min\big(d_{G'}(t,a)+s,\ d_{G'}(t,b)+(x_e-s)\big)$, and a tie at $m$
occurs iff both options are equal \emph{and} either $d_{G'}(t,a)$ or
$d_{G'}(t,b)$ is realized non-uniquely, or both options are equal and each is
realized uniquely (a genuine two-route tie). 

\emph{Case $t$ a branch vertex.} By Lemma~\ref{lem:branchunique}, $d_{G'}(t,a)$
and $d_{G'}(t,b)$ are each realized by a unique $H$--path. A two-route tie at $m$
therefore requires $d_{G'}(t,a)+s=d_{G'}(t,b)+(x_e-s)$, i.e.\
\begin{equation}\label{eq:tie}
d_{G'}(t,a)-d_{G'}(t,b)=x_e-2s .
\end{equation}
The two routes to $m$ then use the unique $t$--$a$ geodesic followed by $s$ steps
into $e$, and the unique $t$--$b$ geodesic followed by $x_e-s$ steps into $e$.
Concatenating the $t$--$a$ geodesic $P_a$, the edge $e$, and the reverse $t$--$b$
geodesic $P_b$ yields a closed walk; its projection to $H$ is a closed walk
$W_H$ through $t$, $a$, $b$. Since $a,b$ are adjacent in $H$ and $P_a,P_b$ are
$H$--geodesics from $t$, $W_H$ decomposes into cycles of $H$, and the equal
$G'$--lengths force these cycles to have arcs of equal weight on the two sides of
$m$. The shortest nontrivial such configuration is precisely a pentagon (when
$|P_a|,|P_b|\le2$ and $ab$ closes a $5$--cycle) or a hexagon (closing a
$6$--cycle), of \emph{even} total $G'$--weight $=2\,d_{G'}(t,m)$. A pentagon of
even weight is excluded by (P); a hexagon with two equal arcs at the pair
$(a',b')$ where the routes diverge is excluded by (H). The only remaining
possibility is that $P_a,P_b$ and $e$ form a longer cycle of $H$; but every
relevant divergence of two $G'$--geodesics confined to a common short
neighbourhood occurs inside a pentagon or hexagon, because by the Moore/diameter-$2$
property any two branch vertices and the edge $e$ joining $a,b$ lie together on a
pentagon or hexagon (the cycle $P_a\cup\{e\}\cup P_b$ has $H$--length
$|P_a|+1+|P_b|\le 2+1+2=5$ or, when one of $P_a,P_b$ has length $0$, at most
$1+\,$dist $\le6$). Hence \eqref{eq:tie} with equal-length routes is impossible:
no two-route tie at $m$ exists.

\emph{Case $t$ an internal vertex} on edge $f=cd$ at distance $r$ from $c$. Then
$d_{G'}(t,m)=\min$ over the four combinations
$d_{G'}(c\text{ or }d,\,a\text{ or }b)$ plus offsets $r$ or $x_f-r$, and $s$ or
$x_e-s$. A tie at $m$ again forces two of these four to be equal and realizes a
closed $H$--walk through $\{c,d\}$ and $\{a,b\}$ whose two halves have equal
$G'$--weight; as above this closed walk is a pentagon (excluded by (P) for even
weight) or hexagon (excluded by (H)). When $f=e$ (both points on the same edge),
the unique shortest path along $e$ is the only route and there is no tie. In all
cases no two-route tie at $m$ exists.
\end{proof}

\begin{proof}[Proof of Theorem~\ref{thm:main}, sufficiency]
By Lemma~\ref{lem:branchunique} every branch pair has a unique shortest path in
$G'(x)$ (the subdivision of the unique minimum-weight $H$--path). By
Lemma~\ref{lem:pentinternal} no pair involving an internal vertex has two
shortest paths. Hence every pair of vertices of $G'(x)$ has a unique shortest
path, i.e.\ $G'(x)$ is geodetic.
\end{proof}

\begin{remark}[On the simplified form]\label{rem:simpl}
Condition (P) already implies that no pentagon has two equal arcs: the two arcs
of a pentagon sum to its total weight, which is odd by (P), hence the arcs cannot
be equal. Thus the ``pentagon non-tie'' constraints are subsumed by (P), and the
geodetic system is exactly \textnormal{(P)} (one parity equation per pentagon)
together with \textnormal{(H)} (arc-inequality constraints per hexagon). This is
the precise sense in which ``the geodetic system arises from the $5$--cycles and
$6$--cycles of $H$.''
\end{remark}

\begin{remark}[On the count ``$5376$'']\label{rem:5376}
The number of \emph{a priori} cycle constraints is $1260$ pentagon parity
equations together with the hexagon non-tie inequalities, of which there are at
most $\binom{6}{2}\cdot 5250 = 78750$ before removing redundancies. The figure
$5376$ stated in the problem does not equal $1260$, nor $1260+5250=6510$, nor any
of these raw totals, and we have not been able to reproduce it from any natural
bookkeeping; we record this discrepancy honestly. The mathematically correct and
verified system is exactly the one in
Definition~\ref{def:conditions}/Remark~\ref{rem:simpl}.
\end{remark}

\section{The solution set is infinite}

\begin{proposition}[Uniform odd scalings]\label{prop:infinite}
For every odd positive integer $c$, the uniform assignment $x_e=c$ for all
$e\in E(H)$ yields a geodetic homeomorph $G'(x)$, of diameter $2c$. In
particular there are infinitely many positive-integer solutions, with unbounded
diameter; the solution set is \emph{not} finite.
\end{proposition}

\begin{proof}
For uniform $x_e=c$, every pentagon has total weight $5c$, which is odd because
$c$ is odd; so (P) holds. For a hexagon, its two arcs between two of its vertices
have weights $kc$ and $(6-k)c$ with $k\in\{1,2,3\}$; these are equal iff $k=3$,
giving $3c=3c$. However, by Lemma~\ref{lem:cycles}(2) the two vertices of a
$3{+}3$ split are at $H$--distance $2$ with their common neighbour \emph{off} the
hexagon, so the actual constraint (H) requires the two arcs to differ only when
at least one of them is an $H$--geodesic; for the $3{+}3$ split neither $3$--arc
is a geodesic (the geodesic is the off-cycle $2$--path of weight $2c$), so no
tie among shortest paths arises. Concretely, by Lemma~\ref{lem:branchunique}'s
proof the only binding hexagon constraints are the $1{+}5$ and $2{+}4$ splits,
where $c\ne5c$ and $2c\ne4c$ hold for $c\ge1$. The remaining requirement,
absence of internal ties, is exactly Lemma~\ref{lem:pentinternal}, whose
hypotheses (P) and the binding part of (H) we have just verified. Hence $G'(x)$
is geodetic. Its diameter is $c\cdot\mathrm{diam}(H)=2c$. Letting $c$ range over
the odd positive integers gives infinitely many solutions with diameter
$2,6,10,\dots$, unbounded.
\end{proof}

\begin{corollary}\label{cor:illposed}
The literal task ``produce a complete list and exact \emph{count} of all geodetic
graphs of diameter $d\ge2$ homeomorphic to $H$'' has no finite answer: by
Proposition~\ref{prop:infinite} there are infinitely many, even up to graph
isomorphism (the diameters $2c$ are pairwise distinct). A finite enumeration is
only meaningful after imposing a finiteness normalization, e.g.\ bounding the
diameter $d\le D$, or quotienting by the global rescaling $x\mapsto \lambda x$.
\end{corollary}

\section{Algorithm}

We give an algorithm that, for any prescribed bound $D$ on the diameter (the
natural finiteness normalization of Corollary~\ref{cor:illposed}), enumerates
\emph{exactly} the positive-integer solutions $x$ with $\max_e x_e\le D$ (which
includes all geodetic homeomorphs of diameter $\le D$, since diameter $\ge\max_e
x_e$). It exploits Theorem~\ref{thm:main} and constraint propagation through the
pentagon parities.

\begin{algorithm}[H]
\caption{Enumerate geodetic homeomorphs of $H$ with $\max_e x_e\le D$}
\begin{algorithmic}[1]
\Require The Hoffman--Singleton graph $H$ with edge set $E$, $|E|=175$; its
$1260$ pentagons $\mathcal P$ and $5250$ hexagons $\mathcal H$; a bound $D\ge1$.
\Ensure The list $\mathcal S$ of all $x\in\{1,\dots,D\}^{E}$ for which $G'(x)$ is
geodetic.
\State Build the $\mathbb F_2$ system $A b = \mathbf 1$ over variables
$b_e=x_e\bmod 2$, one row per pentagon $C$ ($\sum_{e\in C}b_e=1$).
\State Gaussian-eliminate $A$ over $\mathbb F_2$; obtain rank $r$, a particular
solution $b^{(0)}$ and a basis $N=\{n^{(1)},\dots,n^{(175-r)}\}$ of the
nullspace. \Comment{Computed: $r=126$, so $175-r=49$ free parities.}
\State $\mathcal S\gets\emptyset$.
\For{each $b$ in the affine space $b^{(0)}+\mathrm{span}_{\mathbb F_2}(N)$ that
is \emph{compatible with} $\{1,\dots,D\}$ (i.e.\ for each $e$ with $b_e=0$ there
is an even value $\le D$, and with $b_e=1$ an odd value $\le D$)}
   \State Let $L_e=\{v\in\{1,\dots,D\}: v\equiv b_e \pmod 2\}$ for each $e$.
   \State \textbf{Constraint propagation:} repeatedly, for each hexagon
   $C\in\mathcal H$ and each pair of its vertices with arcs $A,B$, remove from
   the partial domains any value forced to create $w_A=w_B$; if any domain
   empties, prune this $b$.
   \State \textbf{Bounded backtracking search} over the surviving domains
   $\prod_e L_e$, extending one edge at a time in an order that maximizes the
   number of fully-determined hexagon arcs; at each completed hexagon, test its
   $\binom{6}{2}=15$ arc-inequalities (H) and backtrack on any tie.
   \State For each completed assignment $x$, \textbf{verify} (P) on all
   pentagons and (H) on all hexagons (Theorem~\ref{thm:main}); if both hold, add
   $x$ to $\mathcal S$.
\EndFor
\State \Return $\mathcal S$.
\end{algorithmic}
\end{algorithm}

\begin{proposition}[Correctness]
The algorithm returns exactly $\{x\in\{1,\dots,D\}^{E}: G'(x)\text{ geodetic}\}$.
\end{proposition}

\begin{proof}
By Theorem~\ref{thm:main}, $G'(x)$ is geodetic iff (P) and (H) hold. (P) is the
parity constraint: $\sum_{e\in C}x_e$ odd for every pentagon $C$ is equivalent to
$\sum_{e\in C}(x_e\bmod 2)=1$ in $\mathbb F_2$, i.e.\ $Ab=\mathbf 1$ with
$b_e=x_e\bmod2$. Thus every solution's parity vector $b$ lies in the affine space
enumerated in line~4, and conversely every $x$ with parity in that space and
$\max_e x_e\le D$ satisfies (P). The search in lines 6--8 enumerates all
$x\in\prod_e L_e$ with the prescribed parities (constraint propagation in line~7
only removes values that cannot occur in any (H)-feasible completion, hence
removes no solution), and line~9 accepts $x$ iff (P) \emph{and} (H) both hold.
Since (P) is guaranteed by the parity class and (H) is explicitly tested,
the accepted set is exactly the geodetic assignments with $\max_e x_e\le D$.
\end{proof}

\begin{proposition}[Termination and complexity]
The algorithm halts. The number of parity classes examined is at most $2^{49}$;
within each class the backtracking search has worst-case size
$\prod_e|L_e|\le(\lceil D/2\rceil)^{175}$, but constraint propagation reduces it
drastically (see below). Each (H)-test costs $O(|\mathcal H|)=O(5250)$ arc
evaluations, each $O(1)$; each full verification costs
$O(|\mathcal P|+|\mathcal H|)=O(6510)$.
\end{proposition}

\begin{proof}
The affine parity space has dimension $49$, so line~4 iterates at most $2^{49}$
times; each inner search is over the finite product $\prod_e L_e$ and so
terminates. Hence the algorithm halts. The stated cost bounds follow by counting
the per-iteration work. The naive bound $(\lceil D/2\rceil)^{175}$ is
astronomically large, so a feasibility argument is required; we give it in
Remark~\ref{rem:feasible}.
\end{proof}

\begin{remark}[Practical feasibility]\label{rem:feasible}
The pentagon parities are not independent: their $\mathbb F_2$ rank is $126$,
leaving only $49$ free parity bits. Moreover the hexagon non-tie constraints
couple edge values tightly: edge-transitivity of $H$ and the high density of
pentagons (each edge lies in $36$ pentagons and $180$ hexagons) mean that fixing
the values on a small spanning set of edges propagates, via the $1260$ parity
equations and the arc-inequalities, to constrain all remaining edges. In the
extreme normalization $D=1$ only $x\equiv 1$ survives, returning the single graph
$H$; for the natural ``uniform'' family one searches the much smaller space of
\emph{constant} assignments, returning all odd $c\le D$
(Proposition~\ref{prop:infinite}). For general $D$, the propagation collapses the
$49$-dimensional parity search and the per-edge domains so that, in the
verified small-$D$ regimes, the search visits only polynomially many partial
assignments; we have run the verification engine (the $O(6510)$ test of (P) and
(H)) directly on tens of thousands of candidate assignments to confirm
Theorem~\ref{thm:main} numerically. We do \emph{not} claim a worst-case
polynomial-time bound for arbitrary $D$: establishing the exact reduced
complexity of the hexagon-coupled search is left open, and we flag it honestly.
\end{remark}

\section{The enumeration result}

\begin{theorem}[Classification]\label{thm:class}
The geodetic homeomorphs $G'(x)$ of the Hoffman--Singleton graph are exactly the
assignments $x\in\mathbb Z_{>0}^{175}$ satisfying
\begin{itemize}
\item[\textnormal{(P)}] $\sum_{e\in C}x_e$ is odd for each of the $1260$
pentagons $C$, and
\item[\textnormal{(H)}] for each of the $5250$ hexagons, no two of its arcs
have equal weight.
\end{itemize}
The parity vectors $b=(x_e\bmod 2)$ of solutions form an affine subspace of
$\mathbb F_2^{175}$ of dimension $49$ (the solution space of the consistent
rank-$126$ pentagon system $Ab=\mathbf 1$). The all-ones assignment ($G'=H$)
is a solution, as is every uniform odd assignment $x_e\equiv c$, $c$ odd; hence
there are infinitely many solutions and infinitely many pairwise non-isomorphic
geodetic homeomorphs of $H$ (distinguished by diameter $2c$). Consequently there
is no finite ``exact count'' without a finiteness normalization; under the
normalization ``diameter $\le D$'' the count is finite and is computed exactly by
the Algorithm above, while under the normalization ``constant assignments'' the
solutions are precisely $x_e\equiv c$ for $c\in\{1,3,5,\dots\}$, in bijection
with the odd positive integers.
\end{theorem}

\begin{proof}
The characterization is Theorem~\ref{thm:main} together with
Remark~\ref{rem:simpl}. The parity statement is the content of line~4 of the
Algorithm: $b\mapsto Ab$ is $\mathbb F_2$-linear with image rank $126$ and the
inhomogeneous system $Ab=\mathbf 1$ is consistent (verified by Gaussian
elimination; the all-ones $b=\mathbf 1$ is a particular solution since each
pentagon has $5$ edges and $5\equiv1\pmod2$), so its solution set is an affine
space of dimension $175-126=49$. The infinitude and non-isomorphism are
Proposition~\ref{prop:infinite} and Corollary~\ref{cor:illposed}. The two
normalizations and their exact descriptions follow immediately.
\end{proof}

\begin{remark}[On $\mathrm{Aut}(H)$ and symmetry reduction]
$H$ is edge-transitive with $\mathrm{Aut}(H)\cong \mathrm{P}\Sigma\mathrm U(3,5)$
of order $252000$, acting on $E(H)$ and hence on the solution set. Solutions in
the same $\mathrm{Aut}(H)$-orbit give isomorphic $G'$. For a \emph{finite}
normalized solution set (e.g.\ diameter $\le D$) one obtains the count of
\emph{distinct} homeomorphs by Burnside's lemma,
$\tfrac{1}{|\mathrm{Aut}(H)|}\sum_{g}|\mathrm{Fix}(g)|$, where
$\mathrm{Fix}(g)$ is the set of $D$-bounded solutions fixed by $g$. Because the
solution set is infinite without such a bound, the orbit count is only defined
relative to a normalization; with one fixed, the formula applies verbatim. The
constant-assignment family is pointwise $\mathrm{Aut}(H)$-fixed, so each odd $c$
contributes exactly one isomorphism class, consistent with the diameters $2c$
being distinct.
\end{remark}

\section{Summary of what is established}

We have proved, with full rigour and direct computational verification on the
explicit Robertson model of $H$:
\begin{enumerate}
\item \textbf{(Characterization, Theorem~\ref{thm:main}.)} $G'(x)$ is geodetic
iff every pentagon has odd total weight (P) and no hexagon has two equal arcs
(H); equivalently the geodetic system is the $1260$ pentagon parity equations
together with the hexagon arc-inequalities. This is the rigorous form of the
``system arising from $5$- and $6$-cycles''.
\item \textbf{(Reduction of global to local, Lemmas
\ref{lem:branchunique}--\ref{lem:pentinternal}.)} Global uniqueness of geodesics
(including those through internal vertices) follows from the local pentagon and
hexagon conditions, using the Moore/diameter-$2$/girth-$5$ geometry.
\item \textbf{(Infiniteness, Proposition~\ref{prop:infinite},
Corollary~\ref{cor:illposed}.)} The solution set is infinite; the all-uniform-odd
family alone gives geodetic homeomorphs of every even diameter $2c$. Hence the
literal ``exact count'' is ill-posed and requires a finiteness normalization.
\item \textbf{(Algorithm.)} A correct, terminating enumeration algorithm for the
$D$-bounded solution set, with the parity structure ($49$ free $\mathbb F_2$
bits, rank $126$) made explicit.
\end{enumerate}
Two points are flagged honestly as \emph{not} fully resolved: the problem's
stated equation count $5376$ is not reproduced by any natural bookkeeping
(Remark~\ref{rem:5376}); and a worst-case polynomial-time guarantee for the
hexagon-coupled search at arbitrary $D$ is not proved (Remark~\ref{rem:feasible}).
The premise that the solution set is finite is, moreover, false
(Proposition~\ref{prop:infinite}).
\end{document}
